% 数学ノート集・共通テンプレート v2（積分論2026年度第1部の版面を参考）
% A5判。章内で定義・定理・例・注意等の番号を共通化。
% コンパイル： platex -interaction=nonstopmode math_note_template_v2.tex (2回)
%              dvipdfmx math_note_template_v2.dvi
\documentclass[a5paper,10pt,onecolumn]{jsarticle}
\usepackage[dvipdfmx,top=22mm,bottom=20mm,inner=19mm,outer=18mm,headheight=13pt,headsep=6mm]{geometry}
\usepackage{amsmath,amssymb,amsthm,mathtools}
\usepackage{enumitem,fancyhdr,titlesec}
\usepackage[dvipdfmx,bookmarks=true,colorlinks=false]{hyperref}
\usepackage{pxjahyper}

% ----- 差し替える文書情報 -----
\newcommand{\NoteTitle}{Laplace変換と反転公式}
\newcommand{\NoteSeries}{数学ノート集}

% ----- 文字・組版 -----
% 参考PDFの質素な明朝＋Computer Modernの印象を引き継ぎながら，
% A5判で窮屈にならないよう本文の行間を少し確保する．
\linespread{1.13}
\setlength{\parindent}{1zw}
\setlength{\parskip}{0pt}
\setlength{\abovedisplayskip}{6pt plus 2pt minus 2pt}
\setlength{\belowdisplayskip}{6pt plus 2pt minus 2pt}
\setlength{\abovedisplayshortskip}{3pt plus 2pt minus 1pt}
\setlength{\belowdisplayshortskip}{4pt plus 2pt minus 1pt}
\setlist[enumerate]{topsep=3pt,itemsep=1pt,parsep=0pt,leftmargin=2.4zw}
\titleformat{\section}{\normalfont\normalsize\bfseries}{\thesection}{0.8zw}{}
\titlespacing*{\section}{0pt}{2.0ex plus .4ex}{.9ex}
\titleformat{\subsection}{\normalfont\normalsize\bfseries}{\thesubsection}{0.7zw}{}
\titlespacing*{\subsection}{0pt}{1.5ex plus .3ex}{.6ex}
\raggedbottom

% ----- 柱とページ番号：ページ上部、罫線なし -----
\fancypagestyle{mathnote}{%
  \fancyhf{}%
  \fancyhead[L]{\footnotesize \NoteSeries\ （\NoteTitle）}%
  \fancyhead[R]{\footnotesize \thepage}%
  \renewcommand{\headrulewidth}{0pt}%
  \renewcommand{\footrulewidth}{0pt}%
}
\pagestyle{mathnote}

% ----- 節内通し番号の定義・定理群 -----
\newtheoremstyle{mathnotetheorem}
 {4.5pt}{4.5pt}{\normalfont}{}%
 {\bfseries}{}{0.65zw}%
 {\thmname{#1}\thmnumber{ #2}\thmnote{（#3）}}
\theoremstyle{mathnotetheorem}
\newtheorem{theorem}{定理}[section]
\newtheorem{definition}[theorem]{定義}
\newtheorem{lemma}[theorem]{補題}
\newtheorem{proposition}[theorem]{命題}
\newtheorem{corollary}[theorem]{系}
\newtheorem{example}[theorem]{例}
\newtheorem{remark}[theorem]{注意}
\renewcommand{\proofname}{証明}
\renewcommand{\qedsymbol}{$\square$}
\makeatletter
\renewenvironment{proof}[1][\proofname]{%
  \par\pushQED{\qed}\normalfont\topsep4\p@\@plus2\p@\relax
  \trivlist\item[\hskip\labelsep\bfseries #1]\ignorespaces
}{\popQED\endtrivlist\@endpefalse}
\makeatother

% ----- 使うノートで適宜拡張する記号 -----
\newcommand{\R}{\mathbb{R}}
\newcommand{\norm}[1]{\lVert #1\rVert}
\newcommand{\abs}[1]{\lvert #1\rvert}

\begin{document}
\section*{読了後にできること}
本稿を読むための前提は，Laplace変換，部分積分，指数関数．
本稿の具体的な到達課題は，Laplace変換を直接計算し，常微分方程式の解に適用する．

\section{Laplace変換の定義}

本稿では $f:[0,\infty)\to\mathbb C$ を可測関数とする．
等式は，特に断らない限りLebesgue積分の意味で解釈する．

\begin{definition}
$f$ に対し，
\[
 \mathcal Lf(s)=F(s)=\int_0^\infty e^{-st}f(t)\,dt
\]
が絶対収束する複素数 $s$ における積分を
\emph{片側Laplace変換}という．
ある実数 $a$ に対して
$e^{-at}f(t)\in L^1(0,\infty)$ となるとき，
$f$ を\emph{指数重み付き可積分}という．
\end{definition}

\begin{proposition}[収束半平面]\label{prop:half-plane}
$e^{-at}f\in L^1(0,\infty)$ と仮定する．
このとき $\mathcal Lf(s)$ は
$\operatorname{Re}s\ge a$ で絶対収束し，
\[
 |F(s)|\le\int_0^\infty e^{-at}|f(t)|\,dt
\]
を満たす．
特に $\operatorname{Re}s>a$ では連続である．
\end{proposition}
\begin{proof}
$\sigma=\operatorname{Re}s\ge a$ なら
$|e^{-st}f(t)|=e^{-\sigma t}|f(t)|
\le e^{-at}|f(t)|$ であるから，
比較によって絶対収束と評価が従う．
$s_n\to s$ に対し被積分関数は各 $t$ で収束し，
上の可積分関数で抑えられるので
Lebesgueの収束定理から $F(s_n)\to F(s)$ となる．
\end{proof}

\begin{example}
$\operatorname{Re}s>\operatorname{Re}\lambda$ に対し
\[
 \mathcal L(e^{\lambda t})(s)=\frac1{s-\lambda}.
\]
また $\operatorname{Re}s>0$ で
$\mathcal L1(s)=1/s$ が成り立つ．
\end{example}
\begin{proof}
$\operatorname{Re}(s-\lambda)>0$ のもとで
\[
 \int_0^\infty e^{-(s-\lambda)t}\,dt
 =\lim_{R\to\infty}
 \frac{1-e^{-(s-\lambda)R}}{s-\lambda}
 =\frac1{s-\lambda}.
\]
後半は $\lambda=0$ とすればよい．
\end{proof}

\begin{definition}
$f:\mathbb R\to\mathbb C$ が可測で
$e^{-\sigma t}f(t)\in L^1(\mathbb R)$ となる
$\sigma\in\mathbb R$ に対して
\[
 \mathcal L_{\mathrm b}f(s)
 =\int_{\mathbb R}e^{-st}f(t)\,dt
 \quad(\operatorname{Re}s=\sigma)
\]
を\emph{両側Laplace変換}という．
\end{definition}

\section{正則性と微分公式}

\begin{theorem}[Laplace変換の正則性]\label{thm:holomorphic}
ある $a\in\mathbb R$ について
$e^{-at}f\in L^1(0,\infty)$ とする．
このとき $F(s)=\mathcal Lf(s)$ は
半平面 $\operatorname{Re}s>a$ で正則であり，
任意の整数 $n\ge0$ に対して
\[
 F^{(n)}(s)
 =(-1)^n\int_0^\infty t^n e^{-st}f(t)\,dt
\]
が成り立つ．
\end{theorem}
\begin{proof}
$\operatorname{Re}s>a$ を固定し，
$0<\delta<\operatorname{Re}s-a$ を取る．
$|h|<\delta/2$ なら
\[
 \frac{e^{-(s+h)t}-e^{-st}}h
 =-t e^{-st}\int_0^1e^{-\theta ht}\,d\theta
\]
である．
絶対値は
$t e^{-(\operatorname{Re}s-\delta/2)t}|f(t)|$
以下であり，
\[
 t e^{-(\operatorname{Re}s-\delta/2-a)t}
\]
は $[0,\infty)$ 上有界なので，
これは定数倍の $e^{-at}|f(t)|$ で抑えられる．
従って優収束定理を用いて
\[
 F'(s)=-\int_0^\infty t e^{-st}f(t)\,dt
\]
を得る．
同じ議論を $t^nf(t)$ に繰り返し適用する．
任意の $n$ と $\varepsilon>0$ に対して
$t^n e^{-\varepsilon t}$ が有界であるから，
必要な可積分支配が常に成り立つ．
従って全ての階数の公式が得られる．
\end{proof}

\begin{proposition}[指数移動と時間平行移動]\label{prop:shifts}
$f$ が指数重み付き可積分とし，
$c\in\mathbb C$，$b\ge0$ とする．
\[
 \mathcal L(e^{ct}f(t))(s)=F(s-c)
\]
が収束領域で成り立つ．
また
\[
 g_b(t)=
 \begin{cases}
 0&(0\le t<b),\\
 f(t-b)&(t\ge b)
 \end{cases}
\]
とすると
\[
 \mathcal Lg_b(s)=e^{-bs}F(s)
\]
が成り立つ．
\end{proposition}
\begin{proof}
第一式は
$e^{-st}e^{ct}=e^{-(s-c)t}$ であることによる．
第二式は変数変換 $u=t-b$ によって
\[
 \int_b^\infty e^{-st}f(t-b)\,dt
 =e^{-bs}\int_0^\infty e^{-su}f(u)\,du
\]
となることから従う．
\end{proof}

\section{時間微分と積分}

\begin{theorem}[微分のLaplace変換]\label{thm:derivatives}
$f:[0,\infty)\to\mathbb C$ が各有限区間で絶対連続で，
$f$ と $f'$ がある $a\in\mathbb R$ に対して
指数重み付き可積分であるとする．
さらに $b>a$ を固定する．
このとき $\operatorname{Re}s>b$ で
\[
 \mathcal L(f')(s)=sF(s)-f(0)
\]
が成り立つ．
同様の仮定を $f,f',\dots,f^{(n)}$ に課せば
\[
 \mathcal L(f^{(n)})(s)
 =s^nF(s)-\sum_{k=0}^{n-1}
 s^{n-1-k}f^{(k)}(0)
\]
を得る．
\end{theorem}
\begin{proof}
$f$ は局所絶対連続なので
$e^{-st}f(t)$ は各有限区間で絶対連続である．
部分積分によって
\[
 \int_0^R e^{-st}f'(t)\,dt
 =e^{-sR}f(R)-f(0)
 +s\int_0^R e^{-st}f(t)\,dt
\]
が成り立つ．
$e^{-a t}f$ と $e^{-a t}f'$ が $L^1$ に属するので
$g(t)=e^{-at}f(t)$ は局所絶対連続であり，
\[
 g'(t)=e^{-at}(f'(t)-af(t))\in L^1(0,\infty).
\]
従って $g$ は $t\to\infty$ で有限の極限を持つ：
$g(t)=g(0)+\int_0^t g'(u)\,du$ による．
特に $g$ は有界であり，
$\operatorname{Re}s>a$ なら
$e^{-sR}f(R)=e^{-(s-a)R}g(R)\to0$ である．
$R\to\infty$ として最初の公式を得る．
高階微分の場合は同じ公式を順に適用し，
帰納法によって表示式を得る．
\end{proof}

\begin{proposition}[積分と時間乗法]\label{prop:integration}
$a\in\mathbb R$ とし，$e^{-at}f(t)\in L^1(0,\infty)$ と仮定して
\[
 H(t)=\int_0^t f(u)\,du
\]
と置く．$\operatorname{Re}s>\max\{a,0\}$ なら
\[
 \mathcal LH(s)=\frac{F(s)}s.
\]
また $\operatorname{Re}s>a$ で
\[
 \mathcal L(tf(t))(s)=-F'(s).
\]
\end{proposition}
\begin{proof}
$\operatorname{Re}s=\sigma>\max\{a,0\}$ とし，
Fubiniの定理を適用する．その正当性は
\[
 \int_0^\infty\int_0^t
 e^{-\sigma t}|f(u)|\,du\,dt
 =\frac1\sigma\int_0^\infty e^{-\sigma u}|f(u)|\,du<\infty
\]
による．従って
\[
 \mathcal LH(s)
 =\int_0^\infty f(u)\int_u^\infty e^{-st}\,dt\,du
 =\frac1s\int_0^\infty e^{-su}f(u)\,du.
\]
後半は定理\ref{thm:holomorphic}で $n=1$ とすればよい．
\end{proof}

\section{畳み込み}

\begin{definition}
$f,g\in L^1_{\mathrm{loc}}(0,\infty)$ に対し
\[
 (f*g)(t)=\int_0^t f(u)g(t-u)\,du
\]
を片側畳み込みという．
\end{definition}

\begin{theorem}[畳み込み定理]\label{thm:convolution}
$f,g$ が
$e^{-at}f,e^{-at}g\in L^1(0,\infty)$ を満たすとする．
このとき $f*g$ はほとんど至る所定義され，
$\operatorname{Re}s\ge a$ において
\[
 \mathcal L(f*g)(s)
 =(\mathcal Lf)(s)(\mathcal Lg)(s)
\]
が成り立つ．
\end{theorem}
\begin{proof}
$t=u+v$ という変数変換に対しTonelliの定理を適用すると
\begin{align*}
 &\int_0^\infty e^{-at}\int_0^t
 |f(u)g(t-u)|\,du\,dt\\
 &\qquad=\int_0^\infty\int_0^\infty
 e^{-a(u+v)}|f(u)||g(v)|\,du\,dv\\
 &\qquad=\norm{e^{-a\cdot}f}_{L^1}
 \norm{e^{-a\cdot}g}_{L^1}<\infty.
\end{align*}
従って畳み込みはほとんど至る所定義され，
$e^{-at}(f*g)(t)$ は可積分である．
$\operatorname{Re}s\ge a$ では同じ絶対値評価が成り立つため，
Fubiniの定理を用いて
\begin{align*}
 \mathcal L(f*g)(s)
 &=\int_0^\infty\int_0^\infty
 e^{-s(u+v)}f(u)g(v)\,du\,dv\\
 &=\left(\int_0^\infty e^{-su}f(u)\,du\right)
   \left(\int_0^\infty e^{-sv}g(v)\,dv\right).
\end{align*}
よって結論が従う．
\end{proof}

\begin{example}
$f(t)=e^{-\alpha t}$，$g(t)=e^{-\beta t}$ なら，
$\alpha\ne\beta$ のとき
\[
 (f*g)(t)=
 \frac{e^{-\alpha t}-e^{-\beta t}}{\beta-\alpha}.
\]
\end{example}
\begin{proof}
\[
 (f*g)(t)=\int_0^t
 e^{-\alpha u}e^{-\beta(t-u)}\,du
 =e^{-\beta t}
 \frac{e^{(\beta-\alpha)t}-1}{\beta-\alpha}
\]
であり，これを整理すればよい．
\end{proof}

\section{Fourier変換とLaplace反転}

\begin{definition}
$f$ が $e^{-ct}f(t)\in L^1(0,\infty)$ を満たすとする．
$g_c:\mathbb R\to\mathbb C$ を
\[
 g_c(t)=\begin{cases}
 e^{-ct}f(t)&(t\ge0),\\
 0&(t<0)
 \end{cases}
\]
と定める．
\end{definition}

\begin{lemma}[垂直直線上のFourier表現]\label{lem:vertical}
$f$ と $c$ が上の仮定を満たすなら
\[
 F(c+i\xi)
 =\int_{\mathbb R}g_c(t)e^{-i\xi t}\,dt
 =\sqrt{2\pi}\,\widehat{g_c}(\xi)
\]
である．ここで右辺のFourier変換は
$(2\pi)^{-1/2}$ 正規化を用いる．
\end{lemma}
\begin{proof}
$g_c(t)=0$（$t<0$）を用いて積分区間を
$[0,\infty)$ に制限すれば
\[
 \int_{\mathbb R}g_c(t)e^{-i\xi t}\,dt
 =\int_0^\infty e^{-(c+i\xi)t}f(t)\,dt
 =F(c+i\xi).
\]
Fourier変換の定義から最後の等式も従う．
\end{proof}

\begin{theorem}[Laplace反転公式]\label{thm:inversion}
$f$ は $e^{-ct}f\in L^1(0,\infty)$ を満たすとする．
さらに $\xi\mapsto F(c+i\xi)$ が $L^1(\mathbb R)$ に属し，
$f$ が $t_0>0$ で連続であると仮定する．
このとき
\[
 f(t_0)=\frac{e^{ct_0}}{2\pi}
 \int_{-\infty}^\infty
 F(c+i\xi)e^{i\xi t_0}\,d\xi.
\]
特に，$f$ が $(0,\infty)$ 上で連続であれば，
この反転公式は各 $t_0>0$ で成立し，Laplace変換は
$f$ を点ごとに一意に決定する．一般の可測函数は
後述の一意性定理により，ほとんど至る所でのみ決定される．
\end{theorem}
\begin{proof}
補題\ref{lem:vertical}によって
$g_c\in L^1(\mathbb R)$ かつ
$\widehat{g_c}\in L^1(\mathbb R)$ である．
Fourier反転定理によって，
$g_c$ の連続点 $t_0$ では
\[
 g_c(t_0)
 =\frac1{2\pi}\int_{\mathbb R}
 F(c+i\xi)e^{i\xi t_0}\,d\xi
\]
が成り立つ．
$g_c(t_0)=e^{-ct_0}f(t_0)$ だから両辺に
$e^{ct_0}$ を掛ければ反転公式となる．
二つの連続函数が同じLaplace変換を持ち，
ともに定理の可積分性の仮定を満たすなら，
各 $t_0>0$ で反転公式を適用できるので，
両者は $(0,\infty)$ 上で一致する．
\end{proof}

\begin{remark}[零集合上の値は復元できない]
可測函数に対する点ごとの一意性は一般には成り立たない．
例えば $f(t)=0$ とし，$g(1)=1$，$g(t)=0$（$t\ne1$）
と定めれば，$f,g$ はすべての $c\in\mathbb R$ に対して
$e^{-ct}f,e^{-ct}g\in L^1(0,\infty)$ を満たし，
両者のLaplace変換は恒等的に零であるが，$f(1)\ne g(1)$．
従って可測函数のLaplace変換は，函数の
ほとんど至る所での同値類のみを決定する．
\end{remark}

\begin{theorem}[Gaussianによる正則化反転]
\label{thm:regularized-inversion}
$e^{-ct}f\in L^1(0,\infty)$ とし，
$f$ は $t_0>0$ で連続とする．
このとき
\[
 f(t_0)=\lim_{\varepsilon\downarrow0}
 \frac{e^{ct_0}}{2\pi}\int_{\mathbb R}
 e^{-\varepsilon\xi^2/2}
 F(c+i\xi)e^{i\xi t_0}\,d\xi.
\]
この式では $\xi\mapsto F(c+i\xi)$ の
絶対可積分性を仮定しない．
\end{theorem}
\begin{proof}
$g_c$ を先の定義の関数とし，
\[
 \rho_\varepsilon(x)
 =(2\pi\varepsilon)^{-1/2}
 e^{-x^2/(2\varepsilon)}
\]
と置く．
GaussianのFourier変換の公式と
Fubiniの定理によって
\[
 (\rho_\varepsilon*g_c)(t_0)
 =\frac1{2\pi}\int_{\mathbb R}
 e^{-\varepsilon\xi^2/2}
 F(c+i\xi)e^{i\xi t_0}\,d\xi.
\]
実際 $|F(c+i\xi)|\le\norm{g_c}_1$ であり，
Gaussian因子が可積分なので全ての積分交換は
絶対収束によって正当化される．
一方，$(\rho_\varepsilon)$ は近似恒等作用素であり，
$g_c$ が $t_0$ で連続で $L^1$ に属することから
\[
 (\rho_\varepsilon*g_c)(t_0)\longrightarrow g_c(t_0)
 =e^{-ct_0}f(t_0).
\]
これを代入して結論を得る．
\end{proof}

\begin{theorem}[Laplace変換の一意性]\label{thm:uniqueness}
$f,g$ が可測で，ある $c\in\mathbb R$ に対して
$e^{-ct}f,e^{-ct}g\in L^1(0,\infty)$ とする．
半平面 $\operatorname{Re}s>c$ において
$\mathcal Lf(s)=\mathcal Lg(s)$ なら
$f=g$ がほとんど至る所成立する．
\end{theorem}
\begin{proof}
$h=f-g$ と置く．
半平面 $\operatorname{Re}s>c$ で $\mathcal Lh=0$ だから，
任意の $b>c$ と $\xi\in\mathbb R$ について
$\mathcal Lh(b+i\xi)=0$ である．
補題\ref{lem:vertical}から
$e^{-bt}h(t)1_{\{t\ge0\}}$ のFourier変換は恒等的に零となる．
$L^1$ 上のFourier変換の一意性定理から
$e^{-bt}h(t)=0$ がほとんど至る所成立する．
$e^{-bt}$ は全ての $t$ で非零なので
$h=0$ がほとんど至る所成立する．
\end{proof}

\section{常微分方程式への応用}

\begin{proposition}[一次線形方程式]\label{prop:ode1}
$a,y_0\in\mathbb C$，$h:[0,\infty)\to\mathbb C$ は連続で
指数型の増大条件
$|h(t)|\le Me^{bt}$ を満たすとする．
初期値問題
\[
 y'(t)+ay(t)=h(t),\qquad y(0)=y_0
\]
の唯一の $C^1$ 解は
\[
 y(t)=y_0e^{-at}
 +\int_0^t e^{-a(t-u)}h(u)\,du
\]
である．
\end{proposition}
\begin{proof}
表示式の右辺を $y(t)$ と置く．
Leibnizの微分公式から
\begin{align*}
 y'(t)
 &=-ay_0e^{-at}+h(t)
   -a\int_0^te^{-a(t-u)}h(u)\,du\\
 &=h(t)-ay(t).
\end{align*}
$t=0$ で積分は零なので $y(0)=y_0$ となる．
任意の二解の差 $z$ は $z'+az=0$，$z(0)=0$ を満たし，
$(e^{at}z(t))'=0$ なので $z=0$ である．
従って唯一解である．
またLaplace変換の微分公式と畳み込み定理を形式上でなく
実際に適用すれば，
十分大きい $\operatorname{Re}s$ に対し
\[
 (s+a)Y(s)=y_0+H(s),\qquad
 Y(s)=\frac{y_0}{s+a}+\frac{H(s)}{s+a}
\]
となる．
右辺をLaplace変換の基本公式で反転すれば
上の表示式を再び得る．
\end{proof}

\begin{proposition}[二階振動方程式]\label{prop:oscillator}
$\omega>0$ とする．初期値問題
\[
 y''(t)+\omega^2y(t)=0,\qquad
 y(0)=a,\quad y'(0)=b
\]
の唯一の $C^2$ 解は
\[
 y(t)=a\cos(\omega t)+
 \frac b\omega\sin(\omega t)
\]
である．
\end{proposition}
\begin{proof}
微分して代入すると表示式は方程式と初期条件を満たす．
二解の差 $z$ について
$E(t)=|z'(t)|^2+\omega^2|z(t)|^2$ と置くと，
$z''=-\omega^2z$ を用いて
\[
 E'(t)=2\operatorname{Re}
 (z''\overline{z'}+\omega^2z'\overline z)=0
\]
となる．$E(0)=0$ だから $z=z'=0$ である．
Laplace変換の微分公式からも
\[
 (s^2+\omega^2)Y(s)=as+b
\]
を得る．
$\mathcal L(\cos\omega t)=s/(s^2+\omega^2)$，
$\mathcal L(\sin\omega t)=\omega/(s^2+\omega^2)$
（$\operatorname{Re}s>0$）なので，
反転と一意性の定理によって同じ解が得られる．
\end{proof}

\begin{remark}
Laplace変換は時間微分を代数的な乗法に変え，
畳み込みを積へ変える．
ただし反転公式には可積分性や連続性の条件が必要である．
本稿のGaussian正則化反転は，
通常の垂直線上の積分が絶対収束しない場合にも
反転を厳密に意味づける一つの方法である．
\end{remark}

\section{到達課題の完全解答}
\noindent\textbf{振動関数のLaplace変換．}
\begin{proposition}
実数 $b$ と $s>0$ に対し
\[
\mathcal L(\sin bt)(s)
=\int_0^\infty e^{-st}\sin(bt)\,dt
=\frac{b}{s^2+b^2}
\]
である．
\end{proposition}
\begin{proof}
$I=\int_0^\infty e^{-st}\sin(bt)dt$，
$J=\int_0^\infty e^{-st}\cos(bt)dt$ と置く．
部分積分より
$I=(b/s)J$，
$J=1/s-(b/s)I$．
これらを解くと
$(s^2+b^2)I=b$ となる．
$s>0$ のため全ての境界項は無限遠で零になる．
\end{proof}
\begin{proposition}
$y''+b^2y=0$，$y(0)=0$，$y'(0)=1$
の解は $b\ne0$ なら
$y(t)=\sin(bt)/b$ である．
\end{proposition}
\begin{proof}
方程式のLaplace変換を取れば
$(s^2+b^2)Y(s)=1$．
上の公式より逆変換は $\sin(bt)/b$ である．
\end{proof}
\noindent\textbf{到達確認．}
$\mathcal L(\cos bt)(s)$ を求めよ．
\par\noindent\textbf{解答．}
部分積分で得た $J=1/s-(b/s)I$ に
$I=b/(s^2+b^2)$ を代入すると
$J=s/(s^2+b^2)$ である．
\end{document}